\chapter{Introduction}
Hello \cite{greenwade93}
\section{Objectives}
\section{Challenges}
\section{Contributions}

I’ll draft a thesis-style interim-report Introduction (about 1–2 pages) that follows your specified flow, cites the key papers you’ve collected, and keeps scope/claims cautious (interim-friendly).


## 1 Introduction

Refactoring is a routine activity in software development: developers restructure code to improve readability, maintainability, and evolvability while intending not to change externally observable behaviour. Classic definitions emphasise behaviour preservation and describe refactoring as a sequence of small, semantics-preserving transformations (Opdyke, 1992; Fowler, 1999/2018), and subsequent surveys systematise these refactoring operators and the surrounding tool support (Mens & Tourwé, 2004). In practice, refactoring is widely encouraged because poorly structured code becomes an obstacle to enhancement and maintenance, and continual adaptation is necessary to keep evolving systems satisfactory (e.g., the motivation articulated in work building on Lehman’s laws and empirical maintenance observations). At the same time, refactoring remains costly: it consumes developer time, introduces risk, and is frequently intertwined with other work such as feature development and bug fixing rather than occurring as isolated “refactoring-only” episodes (Murphy-Hill et al., 2008/2009; Kim et al., 2012).

Most research and tooling that evaluates refactoring focuses on *outcomes* rather than the *process*. A common paradigm is to treat refactoring success as a change in static quality indicators between a “before” and “after” snapshot: reductions in code smells, improvements in cohesion/coupling, decreases in complexity, or related maintainability proxies. Representative threads include metric-driven detection or prioritisation of candidate refactorings (e.g., using metric deltas to distinguish refactoring-like changes from functional modifications, or to select among refactoring options), as well as refactoring recommendation systems that rank alternatives by their impact on design metrics. More recent work expands the signals used for guidance—incorporating traceability information alongside code metrics when ranking refactoring solutions, and exploring how product/process metrics relate to developers’ motivations to refactor (e.g., studies building on mined refactorings and repository history, and proposals that use LLMs to infer motivations from commit artefacts). In parallel, a substantial line of work has improved the ability to *detect* refactorings from code changes, notably by using AST-based differencing and structured matching to infer refactoring operations between revisions (e.g., RefactoringMiner and related approaches), enabling large-scale empirical analyses of refactoring frequency, types, and context (Tsantalis et al., 2018/2020; Falleri et al., 2014).

However, these perspectives largely evaluate refactoring as a mapping from an initial program state to a final program state. They do not directly address whether the *trajectory* taken to reach the end state was efficient, low-risk, or “good” by any explicit criterion. This is an important omission because developers rarely refactor in a single atomic operation: they work through intermediate states, sometimes detour, sometimes perform rework, and often create temporary degradation (e.g., in readability, modular structure, or test/build stability) before arriving at a final outcome. Two developers can produce the same end state via very different paths—one might take many disruptive steps, oscillate between designs, or introduce avoidable churn; another might progress through a more stable sequence of changes. Existing endpoint-based evaluation cannot distinguish these cases. Meanwhile, research on process realism suggests that refactoring is frequently “flossed” into ongoing work rather than performed as a separated “root canal” activity, increasing the likelihood of mixed-intent steps and noisy traces (Murphy-Hill et al., 2008/2009). This makes the quality of the process—how a developer navigates intermediate states—potentially consequential for both developer productivity and software health during evolution.

This project investigates the following thesis question: **given a known start state and a known end state, can we evaluate the quality of the developer’s refactoring process, and identify points where a better process was available?** The core idea is to explicitly model refactoring as a *process* over states rather than only an endpoint comparison. Concretely, the project adopts a state/action/trace framing: snapshots of the codebase are treated as *states*, and developer edits—including IDE refactorings and manual transformations—are treated as *actions* that produce a trace from the start state to the end state. A central component will be a **process-quality metric** that scores traces (and potentially intermediate steps) according to criteria that will be defined and justified using the literature on code quality metrics, smell detection validity, and process signals (e.g., DECOR-style smell detection; empirical findings on metric reliability and disagreement; and work linking process/product factors to refactoring likelihood and developer intent). The metric is not assumed to be purely structural: it may incorporate multiple terms capturing endpoint improvement, intermediate-state degradation, churn/disruption, and “safety” proxies such as build/test preservation, reflecting known limitations of relying on any single signal.

To move beyond scoring a single observed trace, the project will also explore **counterfactual alternatives**: at selected points along the developer’s trace, generate plausible alternative next steps (or short-horizon paths) and compare their scores to the developer’s choice. This aligns with, but differs from, prior work that synthesises or recommends refactoring sequences to reach a target design. For example, ReSynth synthesises sequences of IDE refactorings consistent with a developer’s partial edits to complete a transformation (Raychev et al., 2013), while Refactoring Navigator suggests refactoring steps to guide a system toward a desired architectural target (Lin et al., 2016). These tools demonstrate the feasibility of generating refactoring sequences and navigating toward goals, yet they are not designed to evaluate the quality of a human-chosen process given a user-defined metric. In this project, counterfactual generation is used specifically to support process evaluation—surfacing “you could have done X here” comparisons—rather than to replace the developer’s work or to optimise solely for reachability of the end state.

The expected outcome is an **offline analysis tool** that takes (i) a repository, (ii) a start and end revision, and (iii) an observed trace (e.g., commits or checkpoints), and produces a process-quality assessment: an overall score, step-level breakdown, and highlighted divergence points with higher-scoring local alternatives. IDE integration is considered an extension rather than a requirement for the interim stage; nevertheless, related work on IDE interaction logging and refactoring activity traces suggests that richer traces are feasible under appropriate privacy constraints, and may later support finer-grained “step” definitions than commit boundaries.

This interim report assumes an initial scope sufficient to make measurable progress without overcommitting to a final design. The initial implementation will likely focus on a single language ecosystem (e.g., Java) to leverage mature parsing, metrics tooling, and refactoring mining (e.g., RefactoringMiner). A “step” will initially be treated as a commit or checkpoint between snapshots; finer step segmentation from IDE events or intra-commit edits is an optional extension. The project does not aim, at this stage, to solve global optimal planning over all refactoring operators or to provide exhaustive coverage of every refactoring type; instead, it prioritises a limited, defensible action representation and local counterfactual generation sufficient to evaluate whether meaningful process critiques can be produced.

The remainder of this report reviews related work on refactoring definitions, outcome-based evaluation, metrics and smell detection, refactoring mining, and sequence generation; it then outlines a concrete project plan, an evaluation methodology for assessing both metric usefulness and counterfactual suggestions, and ethical considerations—particularly around trace collection, privacy, and the potential misuse of process analytics for surveillance rather than developer support.

*(Suggested figure for this section: a simple diagram showing Start → s₁ → … → sₖ → … → End, with a branch at sₖ indicating an alternative action leading to a higher-scoring local path.)*
